\section{Main Results}
\label{sec:lower-bound}

\subsection{Main result}
\label{subsec:lb-main}

We use the problem class of Definition~\ref{def:problem-class}, the family
\(\calC_\theta\) of Definition~\ref{def:algorithm-family}, and the complexity
of Definition~\ref{def:complexity-measure}, with the normalized constants
\(\mathbf C_0:=(1,1,0,1,1)\). Since \(L_{\calA}=0\), every instance in
\(\mathcal I(\mathbf C_0,\Delta)\) has an affine inner mapping.

We now present the iteration complexity of \(\calC_\theta\) under the choice
\(\beta=1\). We choose this setting for two reasons. First, when
\(\beta=1\), SPLALM reduces to Smoothed GDA applied to \(\calD_\rho\)
(Table~\ref{tab:special-cases}), whose iteration complexity is known. Second,
Theorem~\ref{thm:lower-bound} below holds for any arbitrary choice of
\((\alpha,\beta)\in[0,1]^2\), including \(\beta=1\). Stating the upper bound
under the same setting yields a direct tightness guarantee.

\begin{theorem}[Iteration complexity of SPLALM with \(\beta=1\)]
\label{thm:upper-beta1}
There exist \(\bar\gamma>0\), depending only on \(\mathbf C\) and \(\rho\),
and \(\bar\tau>0\), depending only on \(\mathbf C\), \(\rho\), and
\(\eta/\gamma\), such that the following holds. Let \(\gamma\le\bar\gamma\),
\(\eta\le\gamma/2\), \(\tau\le\bar\tau\), and \(s\ge\tau\). Then there exist
\(c>0\), depending only on \(\theta\) and \(\mathbf C\), and \(C>0\),
depending only on \(\theta\), \(\mathbf C\), and \(\Delta\), such that, given
any \(\epsilon\in(0,1]\) and \(\alpha:=c\,\epsilon^2\),
\[
T^{\mathrm{KKT}}_\epsilon(\theta,\alpha,1;I)\le C\,\epsilon^{-4}
\qquad\text{for all } I\in\mathcal I(\mathbf C,\Delta).
\]
\end{theorem}
The proof is given in Appendix~\ref{app:upper-beta1}.

\begin{remark}
\label{rem:upper-beta1}
(i) When \(\beta=1\), SPLALM reduces to Smoothed GDA applied to
\(\calD_\rho\). A result analogous to Theorem~\ref{thm:upper-beta1} has been
obtained in~\cite[Theorem~3.4]{zhang2020smoothed}
and~\cite[Theorem~B.1]{li2026smoothing}. However, the initial gap there
depends on the initialization, the dual domain is a given compact set, and
the stationarity measure is evaluated at \((\bx^{k+1},\blmbd^{k+1})\).
Theorem~\ref{thm:upper-beta1} bounds the initial gap by the constants of
\(\mathcal I(\mathbf C,\Delta)\), keeps the multipliers in the interior of a
ball without any projection, and evaluates the KKT residual at
\((\bx^{k+1},\blmbd^k)\), the point that SPLALM outputs.
(ii) The iterates do not depend on \(s\) when \(\beta=1\). Hence the constants
\(c\) and \(C\) in Theorem~\ref{thm:upper-beta1} can be chosen independently
of \(s\): the choices for \(s=\tau\) apply to every \(s\ge\tau\).
The upper bound
requires only one member of \(\calC_\theta\); for \(\beta=0\), the instance
in Section~\ref{subsec:lb-huber} shows that no upper bound holds when
\(\epsilon(\rho+s)\le c_1\).
\end{remark}

\begin{theorem}[Lower bound uniform in \(\rho\) and \(s\)]
\label{thm:lower-bound}
Fix \(0<\eta\le\gamma\), \(\tau>0\), and \(\Delta>0\). There exist constants
\(c_1=c_1(\gamma,\eta,\tau)>0\) and \(c_2=c_2(\gamma,\tau,\Delta)>0\),
independent of \(\rho\), \(s\), \(\alpha\), and \(\beta\), such that the
following holds. For every \(\epsilon\in(0,c_2)\), every \(\rho>0\) and
\(s\ge\tau\) with \(\epsilon(\rho+s)\le c_1\), and every
\((\alpha,\beta)\in[0,1]^2\), there exists an instance
\(I\in\mathcal I(\mathbf C_0,\Delta)\) on which
\(\mathrm{SPLALM}_{\theta,\alpha,\beta}\) requires at least
\[
    \frac{\Delta\,\epsilon^{-4}}{16384\,\gamma\max\{1,4\tau,4\gamma\tau\}}
\]
iterations before one of its outputs is an \(\epsilon\)-KKT point.
Equivalently, 
\[T^{\mathrm{KKT}}_\epsilon(\calC_\theta,\mathcal I(\mathbf C_0,\Delta))
\ge\Delta\epsilon^{-4}/(16384\,\gamma\max\{1,4\tau,4\gamma\tau\})\].
\end{theorem}
Explicit choices of \(c_1\) and \(c_2\) are given in~\eqref{eq:rot-thresholds}.
% \todo{State the GS version $T^{\mathrm{GS}}_\epsilon$ after the strict-convex perturbation lemma is proved.}
% perturbation lemma is proved.

\begin{theorem}[Iteration complexity of SPLALM with \(\beta=0\)]
\label{thm:upper-beta0}
Fix \(\rho,\gamma,\eta,\tau>0\) satisfying the step-size
conditions~\eqref{eq:beta0-stepsizes} for all sufficiently large \(s\).
There exist constants \(c_s>0\) and \(c_\alpha>0\), depending only on
\(\rho,\gamma,\eta,\tau\) and \(\mathbf C\), and \(C'>0\), depending in
addition on \(\Delta\), such that the following holds. Given any
sufficiently small \(\epsilon>0\), let
\[
s:=\frac{c_s}{\epsilon},\qquad
\alpha:=c_\alpha\epsilon,\qquad \beta:=0,\qquad
\theta(\epsilon):=(\rho,s,\gamma,\eta,\tau).
\]
Then
\[
T^{\mathrm{KKT}}_\epsilon(\theta(\epsilon),\alpha,0;I)
\le C'\epsilon^{-3}
\qquad\text{for all }I\in\mathcal I(\mathbf C,\Delta).
\]
All parameters remain constant within each run. The coefficient \(c_s\)
is chosen sufficiently large, and \(c_\alpha\) sufficiently small, as
specified in Appendix~\ref{app:upper-beta0}.
\end{theorem}
The proof is given in Appendix~\ref{app:upper-beta0}.

\begin{remark}
\label{rem:upper-beta0}
(i) When \(\beta=0\), SPLALM reduces to Perturbed Smoothed GDA applied to
\(\calD_\rho\). A result analogous to Theorem~\ref{thm:upper-beta0} has
been obtained in~\cite[Theorem~3.3]{li2026smoothing}.
Theorem~\ref{thm:upper-beta0} bounds the initial gap by the constants of
\(\mathcal I(\mathbf C,\Delta)\), keeps the multipliers in the interior of
a ball without any projection, and evaluates the KKT residual at
\((\bx^{k+1},\blmbd^k)\). The change in the evaluation point contributes
an additional term \(L_h/(10s)\) in~\eqref{eq:beta0-residual-transfer}.
(ii) The coefficient \(c_s\) must be sufficiently large. For fixed \(\rho\),
the choice \(s=c/\epsilon\) with \(0<c<c_1\) satisfies
\(\epsilon(\rho+s)\le c_1\) for all sufficiently small \(\epsilon\), and
Theorem~\ref{thm:lower-bound} applies.
(iii) The primal center stepsize is \(\alpha=\Theta(\epsilon)\), while the
dual center remains zero. Whether a primal center stepsize independent of
\(\epsilon\) attains the same order is outside the scope of this result.
\end{remark}

\begin{corollary}[Scale of the combined penalty]
\label{cor:effective-penalty}
Fix \(0<\eta\le\gamma\), \(\tau>0\), and \(\Delta>0\), and let \(c_1\)
be the constant in Theorem~\ref{thm:lower-bound}.
\begin{enumerate}[label=(\alph*)]
\item Let \(\rho=\rho(\epsilon)>0\) and \(s=s(\epsilon)\ge\tau\) satisfy
\(\rho+s\le c_1/\epsilon\) for all sufficiently small \(\epsilon\).
Then
\[
T^{\mathrm{KKT}}_\epsilon
\bigl(\calC_{\theta(\epsilon)},\mathcal I(\mathbf C_0,\Delta)\bigr)
=\Omega(\epsilon^{-4}).
\]
\item Fix \(\rho>0\) and let \(\rho,\gamma,\eta,\tau\) satisfy the
conditions of Theorem~\ref{thm:upper-beta0} for the class constants
\(\mathbf C_0\). For \(s=c_s/\epsilon\) with the sufficiently large
coefficient specified there, the choice \(\beta=0\) and
\(\alpha=c_\alpha\epsilon\) gives
\[
T^{\mathrm{KKT}}_\epsilon
\bigl(\calC_{\theta(\epsilon)},\mathcal I(\mathbf C_0,\Delta)\bigr)
=O(\epsilon^{-3}).
\]
\end{enumerate}
Thus an \(O(\epsilon^{-3})\) guarantee over
\(\mathcal I(\mathbf C_0,\Delta)\) requires
\(\rho(\epsilon)+s(\epsilon)>c_1/\epsilon\) for all sufficiently small
\(\epsilon\), and is attained with fixed \(\rho\) and
\(s=c_s/\epsilon\) under the conditions in~(b).
The constants \(c_1\) and \(c_s\) differ; the conclusion identifies the
order \(\epsilon^{-1}\) of the required combined penalty.
\end{corollary}
\begin{proof}
Part~(a) follows from Theorem~\ref{thm:lower-bound}. Part~(b) follows from
Theorem~\ref{thm:upper-beta0}, since the infimum in
Definition~\ref{def:complexity-measure} includes
\((\alpha,\beta)=(c_\alpha\epsilon,0)\). If
\(\rho(\epsilon)+s(\epsilon)\le c_1/\epsilon\) along a sequence of
accuracies tending to zero, the lower bound in
Theorem~\ref{thm:lower-bound} precludes an \(O(\epsilon^{-3})\) guarantee
along that sequence. This proves the necessary condition.
\end{proof}
